\documentclass[11pt]{article}
\usepackage[margin=1in]{geometry}
\usepackage{amsmath,amssymb}
\usepackage{hyperref}
\usepackage[numbers,sort&compress]{natbib}
\usepackage{booktabs}

\newcommand{\code}[1]{\texttt{#1}}

\title{Set Theoretic Estimation, Recast in Cohomology\\
\large A standalone, low-jargon note}
\author{Project \texttt{thejeb}}
\date{\today}

\begin{document}
\maketitle

\begin{abstract}
This note explains, in as little jargon as possible, one idea: the
gap between \emph{local agreement} and a \emph{global solution} in set
theoretic estimation (STE) is exactly the kind of gap that cohomology
was invented to measure.  We define every term as it appears, work one
tiny example by hand, and give a plain dictionary between the STE words
and the sheaf/cohomology words.  Everything asserted as a theorem here
is machine-checked in the project's Lean library and cited by name; the
rest is clearly flagged as intuition.  No prior cohomology is assumed.
\end{abstract}

\section{STE in one paragraph}
Set theoretic estimation \citep{combettes1993} models a problem as a
collection of \emph{property sets}.  Each source of information --- a
sensor, a document, a voice --- contributes a set $S_i$ of the
candidate answers consistent with that source.  The \emph{feasible set}
is the intersection
\[
  F \;=\; \bigcap_i S_i,
\]
the answers consistent with \emph{everyone at once}.  The corpus is
\emph{feasible} when $F \neq \varnothing$ and \emph{infeasible} when
$F = \varnothing$.  That is the whole object of study: does a single
answer satisfy every source, and if so, which ones?

\section{The one idea}
Checking $F \neq \varnothing$ directly means reasoning about all
sources and all variables simultaneously --- expensive.  In practice we
would rather check things \emph{locally}: does each variable, looked at
on its own, have a consistent value?  Do each small group of sources
agree on the piece they share?

The catch is that local agreement need not add up to a global answer.
It is possible for every small piece to be internally consistent, and
for every pair of overlapping pieces to agree where they overlap, and
yet for \emph{no} single global assignment to restrict to all of them
at once.  The pieces fit pairwise but cannot be assembled.

\begin{quote}
\textbf{The recast.}  The obstruction to assembling locally-agreeing
pieces into one global solution is a \emph{cohomology class}.  When it
vanishes, local agreement guarantees a global answer, and STE becomes a
local computation.  When it does not vanish, it measures exactly how
much the problem is genuinely coupled --- how far local reasoning falls
short.
\end{quote}

This is the same phenomenon cohomology measures everywhere it appears:
a consistent choice made on every small patch that provably cannot come
from one choice made on the whole.  We now make each word precise.

\section{The building blocks, defined plainly}
Fix a set $V$ of \emph{variables}.  Each variable $v$ ranges over some
value set $A_v$.

\begin{description}
\item[Global configuration.] One choice of a value for every variable
  at once: a function $f$ with $f(v)\in A_v$ for all $v$.  Write the set
  of all of them as $\prod_v A_v$.
\item[Constraint.] A set $T$ of allowed global configurations,
  $T\subseteq\prod_v A_v$.  In STE, $T$ is the feasible set $F$; the
  recast just insists on viewing $F$ as a subset of the full
  configuration space, so we can talk about restricting it.
\item[Context (scope).] A subset $U\subseteq V$ of variables --- a
  ``window'' we are willing to look at all at once.  A small context is
  cheap; the full set $V$ is the expensive global view.
\item[Cover.] A collection of contexts $\{U_j\}$ that together touch
  every variable (every $v$ lies in some $U_j$).  This is how we chop
  the global problem into local windows.
\item[Local section.] A partial assignment on a context $U_j$ --- a
  choice of value for just the variables in that window --- that
  \emph{could} be extended to a full allowed configuration in $T$.  ``A
  local piece that is not obviously doomed.''
\end{description}

\section{Local agreement versus a global solution}
Two definitions carry the whole idea.

\paragraph{Compatible family (local agreement).}
Pick one local section $s_j$ on each context $U_j$, such that any two of
them \emph{agree on the variables their contexts share}:
\[
  \text{for all }j,k \text{ and every } v\in U_j\cap U_k,\qquad
  s_j(v) = s_k(v).
\]
This is ``everything is locally fine and the overlaps match.'' It is the
cheap, checkable condition.  (In cohomology this overlap-agreement is
called the \emph{cocycle condition}.)

\paragraph{Gluing (a global solution).}
The family \emph{glues} if there is a single global configuration
$f\in T$ whose restriction to each window $U_j$ is exactly $s_j$.  One
answer that explains all the local pieces simultaneously.  This is the
expensive thing we actually want.

\paragraph{The obstruction.}
Every genuine global solution obviously produces a compatible family
(just restrict it to each window).  The question is the converse:

\begin{quote}
\emph{Does every compatible family glue?}
\end{quote}

When the answer is \emph{yes} for a constraint $T$ and a cover, we say
\textbf{$T$ glues over that cover}, or that its \emph{\v{C}ech
obstruction vanishes}.  Local agreement then \emph{is} global solvability.
When the answer is \emph{no}, some compatible family is locally perfect
but globally impossible; the number of such stuck families is the size
of the obstruction --- the cohomology.  A vanishing obstruction is the
statement ``$H^1 = 0$'' in this setting; a nonzero one is a nonzero
$H^1$ class.

\section{A dictionary}
\begin{center}
\begin{tabular}{@{}lll@{}}
\toprule
\textbf{STE / plain word} & \textbf{Sheaf-cohomology word} & \textbf{What it means}\\
\midrule
Variable value space $A_v$ & Stalk / local data & values a coordinate can take\\
Constraint / feasible set $T$ & The (sub)presheaf of sections & the allowed answers\\
Context $U_j$ & Open set / patch & a window of variables\\
Partial answer on $U_j$ & Local section & a value on that window\\
Extendable partial answer & Section of the (sub)sheaf & a locally-valid piece\\
Overlaps match & Cocycle / compatibility & pieces agree where they share\\
One answer fits all pieces & Global section / gluing & a genuine global solution\\
Local agreement $\Rightarrow$ global & Sheaf condition / $H^1=0$ & the obstruction vanishes\\
Stuck locally-fine families & Nonzero $H^1$ class & the coupling obstruction\\
\bottomrule
\end{tabular}
\end{center}

\section{One example, by hand}
Take two variables $x,y$, each Boolean ($A_x=A_y=\{0,1\}$), and the
cover by the two single-variable windows $\{x\}$ and $\{y\}$.  There is
no overlap, so ``overlaps match'' is automatic --- \emph{every}
combination of a local value for $x$ and a local value for $y$ is a
compatible family.  There are $2\times 2 = 4$ of them.

\paragraph{A gluable constraint (uncoupled).}
Let $T$ be the ``rectangle'' $\{0,1\}\times\{0,1\}$: anything goes.  All
four compatible families glue --- each is just a global configuration
read locally.  Obstruction: zero.  More generally any \emph{product}
constraint $\prod_v P_v$ (each variable independently confined to some
allowed set $P_v$, no cross-variable link) glues over \emph{every}
cover.

\paragraph{A non-gluable constraint (coupled).}
Let $T$ be the ``diagonal'' $\{(0,0),(1,1)\}$: the two variables are
forced equal (this is the XOR/equality link).  Locally each variable can
still be $0$ or $1$, so all four compatible families are still there.
But only \emph{two} of them --- $(0,0)$ and $(1,1)$ --- come from a
global configuration in $T$.  The mixed family ``$x\mapsto 0,\;
y\mapsto 1$'' is locally flawless and globally impossible.  Two stuck
families remain: the obstruction is $2\neq 0$.  The coupling between $x$
and $y$ is invisible in any single window and shows up \emph{only} as
this obstruction.

That contrast --- rectangle glues, diagonal does not --- is the entire
mechanism in miniature.

\section{The dichotomy, and why it is useful}
The example is not a coincidence.  In the project's Lean library the two
sides are theorems:
\begin{itemize}
\item \textbf{Uncoupled $\Rightarrow$ gluable.}  Every product
  (``rectangular'') constraint glues over every cover
  (\code{rectangular\_cechVanishesCover}).
\item \textbf{Gluable $\Rightarrow$ uncoupled.}  If the obstruction
  vanishes, the constraint \emph{must} be a rectangle
  (\code{CechVanishes} forces \code{IsRectangle}, in module
  \code{Ste.CechObstruction}).
\item \textbf{The witness.}  The diagonal has obstruction exactly $2$,
  a genuinely nonzero class (\code{Ste.CechObstruction}).
\end{itemize}
So, in this setting, \emph{gluable = rectangular = uncoupled} are one
and the same, and the obstruction is a sharp detector of coupling.

Why it matters for solving problems:
\begin{enumerate}
\item \textbf{Gluable problems are cheap.}  If the obstruction
  vanishes, you never assemble a global solution by search.  You solve
  each window --- in the extreme, each single variable --- and local
  agreement is a \emph{theorem} that the global answer exists.  For the
  ``frame'' corpora the project studies (each source asserts unary
  \code{variable = value} facts), feasibility collapses to a
  one-variable-at-a-time scan: a value is consistent iff no two sources
  put different values on the same variable.  This is the same fact as
  ``pairwise agreement forces global agreement'' --- the Helly-number-2
  property (\code{frame\_hasHelly\_two}, module
  \code{Ste.HellyConsistency}) --- and it is why the project's solver
  certifies large corpora by a linear sweep instead of an all-subsets
  search.
\item \textbf{The obstruction localizes the hard part.}  When it does
  \emph{not} vanish, its support tells you \emph{which} variables are
  genuinely entangled and must be reasoned about jointly.  You pay the
  global price only there.  Substitution-cipher key spaces (Carlson's
  setting, \citealp{carlson2012}) are the coupled extreme: they are not
  rectangular, so the obstruction is essential and local reasoning
  cannot finish the job.
\end{enumerate}

In one line: \textbf{the \v{C}ech obstruction is a computable measure of
how far an STE problem is from being solvable window-by-window}, zero
exactly when it factors into independent pieces.

\section{What is actually machine-checked}
To keep the bright line the project insists on, here is what is a
theorem (checked by Lean's kernel, sorry-free) versus what is intuition.

\emph{Theorems} (names are Lean identifiers on the project's main
branch): the compatible-family / gluing / vanishing definitions
(\code{CompatibleFamily}, \code{GluesCover}, \code{CechVanishesCover} in
\code{Ste.CechCover}); rectangles glue over every cover
(\code{rectangular\_cechVanishesCover}); vanishing forces rectangularity
and the diagonal's obstruction is $2$ (\code{Ste.CechObstruction}); the
support-cover form, that vanishing is equivalent to the constraint being
generated by pieces whose scopes fit the cover
(\code{cechVanishesCover\_iff\_exists\_scopedGeneration} in
\code{Ste.SupportCover}); and the frame Helly-2 collapse
(\code{frame\_hasHelly\_two}).

\emph{Intuition / framing, not claimed as proven here:} the word
``cohomology'' as the \emph{full} \v{C}ech $H^1$ machinery --- a cochain
complex with a coboundary map and $H^1 = \ker/\operatorname{im}$ over a
coefficient module --- is the \citet{abramsky2011sheaf,abramsky2012cohomology}
sheaf-theoretic framework the project draws on.  This note uses the
$H^1 = 0$ / obstruction-vanishing \emph{level}, which is what is
mechanized; the graded cochain-complex construction in full generality
is cited, not reproved here.

\section*{The recast, restated}
STE asks whether many sources share a common answer.  Split the question
into windows.  ``Every window is fine and the overlaps match'' is cheap
to check.  Whether that local agreement assembles into one global answer
is governed by a single obstruction: it vanishes exactly for the
uncoupled (rectangular) problems, and then STE is a local computation;
it is nonzero exactly when the problem is genuinely coupled, and then it
measures and locates the coupling.  That obstruction is the cohomology,
and that equivalence is the recast.

\bibliographystyle{plainnat}
\bibliography{../../refs}

\end{document}
